% Figure (S8): what spreading and randomizing capture, by field.
% y = a scheme's gain over no funding as a share of the optimal (complete-information
%     gap-rule) allocation's gain; x = screen noise tau (log). Single round, n = 50,
%     funded share one fifth of the population, 2000 populations per point (max MC SE
%     0.002); screen s = K + N(0, tau). Schemes: equal division among the top fifth
%     by screen, equal division among the top two fifths (half-size grants), partial lottery
%     (half of the top two fifths at random, full-size grants), uniform, lottery over all
%     (a fifth of the population at random). Left panel: heavy-tailed capability
%     (alpha_K = 1.3), tight budget (b = 0.1). Right panel: evenly spread capability
%     (alpha_K = 3.5), b = 0.5. Data: fig11-heavy.dat, fig11-even.dat
%     (price_lotteries.py, 2026-09-08).
\begin{figure}[t]
\centering
\begin{tikzpicture}
\begin{axis}[
  name=left,
  width=0.46\textwidth, height=5.6cm,
  axis lines*=left,
  xmode=log,
  xmin=0.3, xmax=10, ymin=0, ymax=1,
  xtick={0.3, 1, 3, 10},
  xticklabels={0.3, 1, 3, 10},
  ytick={0, 0.25, 0.5, 0.75, 1},
  yticklabels={0, 0.25, 0.5, 0.75, 1},
  xlabel={review noise $\tau$},
  ylabel={share of the optimal allocation's gain},
  ylabel style={font=\small}, xlabel style={font=\small},
  tick label style={font=\small},
  axis line style={line width=0.4pt},
  tick style={line width=0.4pt, black},
  title={\small heavy-tailed, tight budget},
  title style={yshift=-2pt},
  clip=false,
]
\addplot[gray, dashed, line width=0.5pt, domain=0.3:10] {0.445};
\addplot[gray, dotted, line width=0.6pt, domain=0.3:10] {0.378};
\addplot[black, line width=0.7pt, mark=*, mark size=1.1pt] table[x=tau, y=ss] {fig11-heavy.dat};
\addplot[black, line width=0.7pt, mark=o, mark size=1.3pt] table[x=tau, y=ws] {fig11-heavy.dat};
\addplot[black, line width=0.7pt, mark=triangle*, mark size=1.4pt] table[x=tau, y=sl] {fig11-heavy.dat};
\node[font=\scriptsize, anchor=south west] at (axis cs:0.315,0.79) {top fifth by review};
\node[font=\scriptsize, gray, anchor=north west] at (axis cs:0.315,0.372) {lottery over all};
\end{axis}
\begin{axis}[
  at={(left.outer east)}, anchor=outer west, xshift=6pt,
  width=0.46\textwidth, height=5.6cm,
  axis lines*=left,
  xmode=log,
  xmin=0.3, xmax=10, ymin=0, ymax=1,
  xtick={0.3, 1, 3, 10},
  xticklabels={0.3, 1, 3, 10},
  ytick={0, 0.25, 0.5, 0.75, 1},
  yticklabels={,,,,},
  xlabel={review noise $\tau$},
  xlabel style={font=\small},
  tick label style={font=\small},
  axis line style={line width=0.4pt},
  tick style={line width=0.4pt, black},
  title={\small evenly spread, ample budget},
  title style={yshift=-2pt},
  clip=false,
]
\addplot[gray, dashed, line width=0.5pt, domain=0.3:10] {0.831};
\addplot[gray, dotted, line width=0.6pt, domain=0.3:10] {0.414};
\addplot[black, line width=0.7pt, mark=*, mark size=1.1pt] table[x=tau, y=ss] {fig11-even.dat};
\addplot[black, line width=0.7pt, mark=o, mark size=1.3pt] table[x=tau, y=ws] {fig11-even.dat};
\addplot[black, line width=0.7pt, mark=triangle*, mark size=1.4pt] table[x=tau, y=sl] {fig11-even.dat};
\node[font=\scriptsize, gray, anchor=south west] at (axis cs:0.32,0.836) {uniform};
\node[font=\scriptsize, anchor=south west] at (axis cs:1.15,0.705) {top two fifths by review};
\end{axis}
\end{tikzpicture}
\caption{What five funding schemes gain, by field. Each scheme's gain over no funding
as a share of the optimal allocation's gain, in a single round with a fifth of the
population funded, as review noise increases. Filled circles: an equal division among
the top fifth of researchers by review score. Open circles: an equal division among
the top two fifths. Triangles: the partial lottery, which funds half of the top two
fifths at random. Dashed rule: uniform funding. Dotted rule: a lottery over all
researchers. In a heavy-tailed field with a tight budget (left), selection by review
gains most, and its gain declines slowly as review grows noisier; the partial lottery
gains less, mainly because it ignores the review ranking within the top two fifths;
the lottery over all researchers gains least, less than uniform funding. In an evenly
spread field with a large budget (right), the ordering reverses: uniform funding gains
most, and every scheme that concentrates grants gains less. Two thousand populations
per point.}
\label{fig:lottery-schemes}
\end{figure}
